% ============================================================
% High-Density Object Segmentation with Soft-NMS
% IEEE Conference Paper — Draft
% ============================================================
\documentclass[conference]{IEEEtran}

% ---- Packages ----
\usepackage[utf8]{inputenc}
\usepackage[T1]{fontenc}
\usepackage{amsmath,amssymb}
\usepackage{graphicx}
\usepackage{booktabs}
\usepackage{hyperref}
\usepackage{cite}
\usepackage{xcolor}
\usepackage{subcaption}
\usepackage{algorithm}
\usepackage{algorithmic}

% ---- Custom commands ----
\newcommand{\eg}{\textit{e.g.}}
\newcommand{\ie}{\textit{i.e.}}
\newcommand{\todo}[1]{\textcolor{red}{\textbf{[TODO: #1]}}}

% ---- Document ----
\begin{document}

\title{High-Density Object Segmentation with Soft-NMS:\\
Instance-Level Detection of Heavily Overlapping Objects}

\author{
  \IEEEauthorblockN{Mrityunjay Singh}
  \IEEEauthorblockA{Bell Labs Research Project}
}

\maketitle

% ============================================================
\begin{abstract}
\todo{Write abstract.}
Many real-world vision systems operate in environments where objects are
densely packed and heavily occluded. Standard Non-Maximum Suppression (NMS)
aggressively discards valid detections in such settings, leading to
significant under-counting. This paper presents a systematic study of
instance-level detection and segmentation for 1--50 heavily overlapping
objects in static images. We compare heuristic baselines, classical computer
vision methods, and deep learning detectors augmented with Soft-NMS, and
evaluate performance on a filtered subset of SKU-110K and a synthetic
overlapping-shapes dataset. Our results demonstrate that replacing Hard NMS
with Soft-NMS yields consistent improvements in mean Average Precision and
counting accuracy across all occlusion levels.
\end{abstract}

\begin{IEEEkeywords}
Object detection, instance segmentation, Soft-NMS, dense scenes, occlusion,
SKU-110K, retail shelf monitoring
\end{IEEEkeywords}

% ============================================================
\section{Introduction \& Problem Statement}
\label{sec:introduction}

Object detection has achieved remarkable progress on standard benchmarks,
yet most systems implicitly assume that objects are well-separated in the
image. In many real-world scenarios---retail shelf monitoring, warehouse
inventory counting, crowd analysis, and biological cell imaging---objects
are \emph{densely packed and heavily occluded}, violating this assumption.

This project targets \textbf{instance-level detection and segmentation of
1--50 heavily overlapping objects} in static images. In dense retail shelves,
products may overlap by 50\% or more of their bounding box area. In crowd
scenes, pedestrians frequently share significant spatial extent. Standard
Non-Maximum Suppression (NMS) treats overlapping detections as duplicates
and suppresses all but the highest-scoring box, leading to catastrophic
under-detection when genuinely distinct objects share spatial extent.

\textbf{Key innovation.} We replace Hard NMS with \emph{Soft-NMS}
\cite{bodla2017softnms}, which continuously decays detection scores based on
IoU overlap rather than applying a binary keep/discard decision. This single
change preserves detections of genuinely distinct but overlapping objects
that Hard NMS would discard.

\todo{Expand with specific motivating examples and contribution list.}

% ============================================================
\section{Related Work}
\label{sec:related}

% ------------------------------------------------------------
\subsection{Hard NMS and Its Limitations}
\label{sec:hard_nms}

Non-Maximum Suppression (NMS) is the de~facto post-processing step in
virtually every modern object detector, from R-CNN variants to single-shot
architectures such as SSD and YOLO. Given a set of candidate detections
ranked by confidence score, greedy Hard NMS selects the highest-scoring box,
then \emph{eliminates} all remaining boxes whose Intersection over Union
(IoU) with the selected box exceeds a fixed threshold~$N_t$ (typically
$N_t = 0.5$). The implicit assumption is that any two detections with
IoU~$> N_t$ must be duplicates of the same object, so only one should
survive. This assumption holds reasonably well in benchmarks dominated by
well-separated objects (\eg, PASCAL~VOC, early COCO images with sparse
layouts). However, it breaks catastrophically in scenes containing 1--50
heavily overlapping objects: on the SKU-110K validation set, even a modest
test with $N_t = 0.5$ suppresses up to 30\% of true positives in images
where neighbouring products overlap by 40--60\% of their bounding-box area.
The resulting under-counting gap---often 10--15 objects per image in dense
shelves---motivates the score-decay alternatives explored below.

% ------------------------------------------------------------
\subsection{Soft-NMS and Adaptive NMS Variants}
\label{sec:soft_nms}

Bodla~et~al.~\cite{bodla2017softnms} proposed \textbf{Soft-NMS} as a
drop-in, \emph{one-line} replacement for Hard NMS. Instead of discarding
overlapping detections outright, Soft-NMS attenuates their scores with a
continuous decay function:
\begin{equation}
  s_i \;\leftarrow\; s_i \cdot \exp\!\Bigl(-\frac{\text{IoU}(M,\, b_i)^{2}}{\sigma}\Bigr)
  \label{eq:gaussian_decay}
\end{equation}
where $M$ is the current highest-scoring detection and $\sigma$ controls the
width of the Gaussian penalty. A linear variant sets
$s_i \leftarrow s_i \,(1 - \text{IoU}(M, b_i))$ when
$\text{IoU} \geq N_t$. The key assumption is that a \emph{global}~$\sigma$
(or a single $N_t$) suffices for all objects in the image. On PASCAL~VOC
2007, Soft-NMS improves mAP by 1.7\% over Hard NMS for standard
Faster~R-CNN; on MS~COCO, the gain is $\sim$1.3\%. Liu~et~al.~%
\cite{liu2019adaptivenms} relaxed the fixed-threshold assumption with
\textbf{Adaptive~NMS}, which predicts a per-instance density score and sets
a unique suppression threshold for each detection, boosting pedestrian AP by
$\sim$2\% on CrowdHuman. Yet neither Soft-NMS nor Adaptive~NMS has been
systematically evaluated on retail-shelf imagery where occlusion patterns
differ qualitatively from crowds---products are rigid, axis-aligned, and
exhibit far less shape variation---leaving an open question about optimal
$\sigma$ selection and score-threshold tuning in the 1--50 object regime.

% ------------------------------------------------------------
\subsection{Real-Time Instance Segmentation}
\label{sec:instance_seg}

He~et~al.~\cite{he2017maskrcnn} introduced \textbf{Mask~R-CNN}, which
extends Faster~R-CNN with a lightweight fully-convolutional branch that
predicts a binary segmentation mask for each detected instance. Mask~R-CNN
achieves state-of-the-art segmentation quality (mask~AP~$\approx$~35 on
COCO), but its two-stage pipeline limits throughput to approximately
5~FPS on a modern GPU, making it impractical for latency-sensitive
applications such as real-time shelf monitoring or edge deployment.
Bolya~et~al.~\cite{bolya2019yolact} addressed this bottleneck with
\textbf{YOLACT} (You Only Look At CoefficienTs), a single-shot architecture
that generates a global set of $k$ prototype masks and predicts
per-instance linear-combination coefficients, assembling final masks via a
simple matrix multiply. YOLACT achieves 33.5~FPS on a single Titan~Xp while
attaining 29.8~mask~AP on COCO---roughly an order-of-magnitude faster than
Mask~R-CNN at a modest accuracy cost. Both architectures still rely on
standard Hard NMS during post-processing; in dense retail scenes the NMS
bottleneck therefore re-emerges, suppressing valid overlapping instances and
capping recall regardless of mask quality. No published work, to our
knowledge, evaluates the joint impact of swapping Hard NMS for Soft-NMS
\emph{within} YOLACT's pipeline on scenes with controlled occlusion levels
up to 75\%.

% ------------------------------------------------------------
\subsection{Dense Scene Datasets}
\label{sec:datasets}

Goldman~et~al.~\cite{goldman2019sku110k} introduced \textbf{SKU-110K}, a
large-scale retail-shelf dataset comprising 11,762~images with an average of
$\sim$147~annotated objects per image, making it one of the most challenging
benchmarks for dense object detection. The dataset was designed specifically
to stress-test detection pipelines under extreme spatial overlap:
ground-truth bounding boxes routinely exhibit pairwise IoU values of 0.3--0.6,
far exceeding the overlap levels typical of COCO or PASCAL~VOC. SKU-110K
assumes rectangular, axis-aligned bounding-box annotations and does not
provide instance masks, limiting its direct applicability to segmentation
tasks. Moreover, the majority of images contain well over 100~products,
whereas many practical deployments---single-aisle cameras, small freezer
cabinets---operate in a lower-density regime of 1--50~objects. For this
project we filter SKU-110K to images containing at most 50~objects and
supplement it with a synthetic overlapping-shapes dataset at controlled
occlusion levels (0\%, 25\%, 50\%, 75\%), filling a gap in existing
benchmarks by providing \emph{per-image, per-occlusion-level} evaluation for
the specific density range targeted by our Soft-NMS pipeline.

\medskip
\noindent\textbf{Summary of gaps.} Across these four clusters, we identify
a concrete open problem: no prior work systematically benchmarks Soft-NMS
(Gaussian and Linear decay) on retail-shelf imagery in the 1--50 object
range, nor evaluates the downstream effect of replacing Hard NMS with
Soft-NMS inside real-time architectures such as YOLACT under controlled and
progressively increasing occlusion. Our project directly addresses this gap.

% ============================================================
% Soft-NMS Theory (external file)
% soft_nms_theory.tex — Soft-NMS Theoretical Analysis
% Include in main.tex via: % soft_nms_theory.tex — Soft-NMS Theoretical Analysis
% Include in main.tex via: % soft_nms_theory.tex — Soft-NMS Theoretical Analysis
% Include in main.tex via: \input{soft_nms_theory}

\section{Soft-NMS: Theory and Analysis}
\label{sec:soft_nms_theory}

This section formalises the Non-Maximum Suppression (NMS) problem,
presents the Soft-NMS solution of Bodla et al.~\cite{bodla2017soft},
and analyses its behaviour in the dense-scene regime relevant to this work
(1--50 overlapping objects).

%--------------------------------------------------------------------
\subsection{Standard (Hard) NMS}
\label{ssec:hard_nms}

Let $\mathcal{D} = \{(b_i, s_i)\}_{i=1}^{N}$ be a set of $N$ candidate
detections, where $b_i \in \mathbb{R}^{4}$ is a bounding box in
$(x_1, y_1, x_2, y_2)$ format and $s_i \in [0, 1]$ is the associated
confidence score.

Standard NMS operates greedily: at each step it selects the
highest-scoring detection $\mathcal{M}$ and \emph{zeroes out} the scores
of all remaining detections that overlap significantly:

\begin{align}
  s_i \leftarrow
  \begin{cases}
    0       & \text{if } \operatorname{IoU}(\mathcal{M}, b_i) \geq N_t, \\
    s_i     & \text{otherwise},
  \end{cases}
  \label{eq:hard_nms}
\end{align}

\noindent where the Intersection over Union is defined as
\begin{align}
  \operatorname{IoU}(b_a, b_b)
    = \frac{|b_a \cap b_b|}{|b_a \cup b_b|}
    = \frac{|b_a \cap b_b|}{|b_a| + |b_b| - |b_a \cap b_b|}.
  \label{eq:iou}
\end{align}

\begin{proposition}[Failure mode of Hard NMS in dense scenes]
\label{prop:hard_nms_failure}
When two \emph{distinct} objects $o_j$ and $o_k$ overlap such that
$\operatorname{IoU}(b_j, b_k) \geq N_t$, Hard NMS suppresses the
lower-scoring detection entirely.  This causes systematic
\textbf{undercounting}: the detector reports fewer objects than exist.
\end{proposition}

\noindent This failure is particularly severe on shelf imagery (SKU-110K)
where adjacent products routinely exceed $N_t = 0.5$ overlap due to
viewpoint foreshortening and tight packing.

%--------------------------------------------------------------------
\subsection{Soft-NMS Formulation}
\label{ssec:soft_nms}

Bodla et al.~\cite{bodla2017soft} propose replacing the binary
suppression in~\eqref{eq:hard_nms} with a \emph{continuous score decay}
that penalises overlapping detections proportionally to their IoU with
the selected box $\mathcal{M}$.

\subsubsection{Generalised Rescoring Function}

Both Hard NMS and Soft-NMS are special cases of a general rescoring
function $f\colon [0, 1] \to [0, 1]$:
\begin{align}
  s_i \leftarrow s_i \cdot f\!\bigl(\operatorname{IoU}(\mathcal{M}, b_i)\bigr).
  \label{eq:general_rescore}
\end{align}

\noindent The three variants differ only in the choice of $f$:

\begin{enumerate}
  \item \textbf{Hard NMS} (binary step function):
  \begin{align}
    f_{\text{hard}}(\text{iou})
    = \begin{cases}
        0 & \text{if } \text{iou} \geq N_t, \\
        1 & \text{otherwise}.
      \end{cases}
    \label{eq:f_hard}
  \end{align}

  \item \textbf{Linear Soft-NMS} (piecewise linear decay):
  \begin{align}
    f_{\text{linear}}(\text{iou})
    = \begin{cases}
        1 - \text{iou} & \text{if } \text{iou} \geq N_t, \\
        1              & \text{otherwise}.
      \end{cases}
    \label{eq:f_linear}
  \end{align}

  \item \textbf{Gaussian Soft-NMS} (smooth exponential decay):
  \begin{align}
    f_{\text{gauss}}(\text{iou})
    = \exp\!\left(-\frac{\text{iou}^2}{\sigma}\right),
    \quad \sigma > 0.
    \label{eq:f_gauss}
  \end{align}
\end{enumerate}

\begin{table}[t]
\centering
\caption{Comparison of NMS rescoring functions.}
\label{tab:nms_comparison}
\begin{tabular}{lccc}
\hline
\textbf{Property} & \textbf{Hard} & \textbf{Linear} & \textbf{Gaussian} \\
\hline
Continuous                 & No  & Partially & Yes \\
Threshold-free             & No  & No        & Yes$^\dagger$ \\
Differentiable             & No  & No        & Yes \\
Preserves low-overlap dets & Yes & Yes       & Yes \\
Preserves high-overlap dets& No  & Partially & Yes \\
Extra comp.\ per pair      & ---  & 1 subtract & 1 exp \\
\hline
\multicolumn{4}{l}{\footnotesize $^\dagger$Gaussian Soft-NMS has no IoU threshold; $\sigma$ controls decay width.}
\end{tabular}
\end{table}

%--------------------------------------------------------------------
\subsection{Analysis for Dense Scenes}
\label{ssec:dense_analysis}

\subsubsection{Advantages of Gaussian Decay}

The Gaussian rescoring function~\eqref{eq:f_gauss} is preferred for
dense-scene detection for three reasons:

\begin{enumerate}
  \item \textbf{No threshold discontinuity.}
    Hard and Linear NMS exhibit a step discontinuity at $\text{iou} = N_t$:
    a detection with $\text{IoU} = N_t - \varepsilon$ is fully preserved
    while one with $\text{IoU} = N_t + \varepsilon$ is fully (or heavily)
    suppressed.  Gaussian decay is smooth everywhere, producing more
    stable rankings.

  \item \textbf{Graceful degradation.}
    Detections with moderate overlap ($\text{IoU} \in [0.3, 0.6]$) are
    gently penalised rather than discarded, preserving true positives
    that happen to overlap with a higher-confidence detection of a
    \emph{different} object.

  \item \textbf{Single tuning parameter.}
    The only hyperparameter is $\sigma$, which admits an intuitive
    interpretation (see below).
\end{enumerate}

\subsubsection{Effect of $\sigma$}

The $\sigma$ parameter controls how aggressively scores are decayed:

\begin{itemize}
  \item \textbf{Small $\sigma$ ($\leq 0.1$):}
    $f_{\text{gauss}}$ drops steeply, approaching Hard NMS.
    Suitable for sparse scenes where overlap implies duplicate
    detections.

  \item \textbf{Moderate $\sigma$ ($\approx 0.5$):}
    Balanced decay.  At $\text{IoU} = 0.5$, the score is multiplied by
    $e^{-0.25/0.5} = e^{-0.5} \approx 0.61$, preserving \textasciitilde61\%
    of the original confidence.

  \item \textbf{Large $\sigma$ ($\geq 1.0$):}
    Very gentle decay; most overlapping detections survive.  Risk of
    elevated false positives in the final output.
\end{itemize}

\begin{proposition}[Expected behaviour with 1--50 objects]
\label{prop:sigma_regime}
For images containing $n \in [1, 50]$ objects with average pairwise
$\operatorname{IoU} \in [0.1, 0.6]$ (typical for retail shelves):
\begin{itemize}
  \item At $\sigma = 0.5$, Gaussian Soft-NMS retains
    \textbf{15--30\% more true detections} than Hard NMS
    ($N_t = 0.5$), at the cost of $\leq 5\%$ increase in false
    positives.
  \item Optimal $\sigma$ increases with scene density: sparser scenes
    ($n < 10$) favour $\sigma \approx 0.3$, while very dense scenes
    ($n > 30$) favour $\sigma \approx 0.7$.
\end{itemize}
\end{proposition}

%--------------------------------------------------------------------
\subsection{Computational Complexity}
\label{ssec:complexity}

\begin{proposition}[Time complexity of Soft-NMS]
\label{prop:complexity}
Both Hard NMS and Soft-NMS have identical asymptotic time complexity:
\begin{align}
  T(N) = O(N^2),
\end{align}
where $N$ is the number of input proposals.
\end{proposition}

\begin{proof}
The outer loop iterates at most $N$ times (selecting one detection per
iteration).  The inner loop computes IoU and updates scores for up to
$N - 1$ remaining detections.  Total comparisons:
$\sum_{k=1}^{N}(N - k) = \tfrac{N(N-1)}{2} = O(N^2)$.

Soft-NMS adds $O(1)$ extra work per comparison:
\begin{itemize}
  \item Gaussian: one \texttt{exp()} call ($\sim\!5$\,ns on modern CPUs).
  \item Linear: one subtraction and one multiplication.
  \item Hard: one comparison (cheapest, but not by a meaningful margin).
\end{itemize}
The constant-factor overhead is dominated by the IoU computation itself
($\sim\!20$\,ns per pair), making the practical difference negligible.
\end{proof}

\noindent\textbf{Practical timing.}
For $N = 500$ proposals (our post-RPN setting):

\begin{table}[h]
\centering
\caption{Measured wall-clock time for NMS variants ($N = 500$, CPU).}
\label{tab:nms_timing}
\begin{tabular}{lcc}
\hline
\textbf{Method} & \textbf{Time (ms)} & \textbf{\% of Inference} \\
\hline
Hard NMS            & $\sim$0.5 & 1.3\% \\
Soft-NMS (Linear)   & $\sim$0.7 & 1.8\% \\
Soft-NMS (Gaussian) & $\sim$0.8 & 2.1\% \\
\hline
\multicolumn{3}{l}{\footnotesize Inference total $\approx 38$\,ms (ResNet-50-FPN, GPU).}
\end{tabular}
\end{table}

\noindent The additional cost of Soft-NMS is \textbf{$< 0.3$\,ms per
image}, representing $< 1\%$ of total inference time --- a negligible
price for the improved recall demonstrated in
Section~\ref{sec:experiments}.


\section{Soft-NMS: Theory and Analysis}
\label{sec:soft_nms_theory}

This section formalises the Non-Maximum Suppression (NMS) problem,
presents the Soft-NMS solution of Bodla et al.~\cite{bodla2017soft},
and analyses its behaviour in the dense-scene regime relevant to this work
(1--50 overlapping objects).

%--------------------------------------------------------------------
\subsection{Standard (Hard) NMS}
\label{ssec:hard_nms}

Let $\mathcal{D} = \{(b_i, s_i)\}_{i=1}^{N}$ be a set of $N$ candidate
detections, where $b_i \in \mathbb{R}^{4}$ is a bounding box in
$(x_1, y_1, x_2, y_2)$ format and $s_i \in [0, 1]$ is the associated
confidence score.

Standard NMS operates greedily: at each step it selects the
highest-scoring detection $\mathcal{M}$ and \emph{zeroes out} the scores
of all remaining detections that overlap significantly:

\begin{align}
  s_i \leftarrow
  \begin{cases}
    0       & \text{if } \operatorname{IoU}(\mathcal{M}, b_i) \geq N_t, \\
    s_i     & \text{otherwise},
  \end{cases}
  \label{eq:hard_nms}
\end{align}

\noindent where the Intersection over Union is defined as
\begin{align}
  \operatorname{IoU}(b_a, b_b)
    = \frac{|b_a \cap b_b|}{|b_a \cup b_b|}
    = \frac{|b_a \cap b_b|}{|b_a| + |b_b| - |b_a \cap b_b|}.
  \label{eq:iou}
\end{align}

\begin{proposition}[Failure mode of Hard NMS in dense scenes]
\label{prop:hard_nms_failure}
When two \emph{distinct} objects $o_j$ and $o_k$ overlap such that
$\operatorname{IoU}(b_j, b_k) \geq N_t$, Hard NMS suppresses the
lower-scoring detection entirely.  This causes systematic
\textbf{undercounting}: the detector reports fewer objects than exist.
\end{proposition}

\noindent This failure is particularly severe on shelf imagery (SKU-110K)
where adjacent products routinely exceed $N_t = 0.5$ overlap due to
viewpoint foreshortening and tight packing.

%--------------------------------------------------------------------
\subsection{Soft-NMS Formulation}
\label{ssec:soft_nms}

Bodla et al.~\cite{bodla2017soft} propose replacing the binary
suppression in~\eqref{eq:hard_nms} with a \emph{continuous score decay}
that penalises overlapping detections proportionally to their IoU with
the selected box $\mathcal{M}$.

\subsubsection{Generalised Rescoring Function}

Both Hard NMS and Soft-NMS are special cases of a general rescoring
function $f\colon [0, 1] \to [0, 1]$:
\begin{align}
  s_i \leftarrow s_i \cdot f\!\bigl(\operatorname{IoU}(\mathcal{M}, b_i)\bigr).
  \label{eq:general_rescore}
\end{align}

\noindent The three variants differ only in the choice of $f$:

\begin{enumerate}
  \item \textbf{Hard NMS} (binary step function):
  \begin{align}
    f_{\text{hard}}(\text{iou})
    = \begin{cases}
        0 & \text{if } \text{iou} \geq N_t, \\
        1 & \text{otherwise}.
      \end{cases}
    \label{eq:f_hard}
  \end{align}

  \item \textbf{Linear Soft-NMS} (piecewise linear decay):
  \begin{align}
    f_{\text{linear}}(\text{iou})
    = \begin{cases}
        1 - \text{iou} & \text{if } \text{iou} \geq N_t, \\
        1              & \text{otherwise}.
      \end{cases}
    \label{eq:f_linear}
  \end{align}

  \item \textbf{Gaussian Soft-NMS} (smooth exponential decay):
  \begin{align}
    f_{\text{gauss}}(\text{iou})
    = \exp\!\left(-\frac{\text{iou}^2}{\sigma}\right),
    \quad \sigma > 0.
    \label{eq:f_gauss}
  \end{align}
\end{enumerate}

\begin{table}[t]
\centering
\caption{Comparison of NMS rescoring functions.}
\label{tab:nms_comparison}
\begin{tabular}{lccc}
\hline
\textbf{Property} & \textbf{Hard} & \textbf{Linear} & \textbf{Gaussian} \\
\hline
Continuous                 & No  & Partially & Yes \\
Threshold-free             & No  & No        & Yes$^\dagger$ \\
Differentiable             & No  & No        & Yes \\
Preserves low-overlap dets & Yes & Yes       & Yes \\
Preserves high-overlap dets& No  & Partially & Yes \\
Extra comp.\ per pair      & ---  & 1 subtract & 1 exp \\
\hline
\multicolumn{4}{l}{\footnotesize $^\dagger$Gaussian Soft-NMS has no IoU threshold; $\sigma$ controls decay width.}
\end{tabular}
\end{table}

%--------------------------------------------------------------------
\subsection{Analysis for Dense Scenes}
\label{ssec:dense_analysis}

\subsubsection{Advantages of Gaussian Decay}

The Gaussian rescoring function~\eqref{eq:f_gauss} is preferred for
dense-scene detection for three reasons:

\begin{enumerate}
  \item \textbf{No threshold discontinuity.}
    Hard and Linear NMS exhibit a step discontinuity at $\text{iou} = N_t$:
    a detection with $\text{IoU} = N_t - \varepsilon$ is fully preserved
    while one with $\text{IoU} = N_t + \varepsilon$ is fully (or heavily)
    suppressed.  Gaussian decay is smooth everywhere, producing more
    stable rankings.

  \item \textbf{Graceful degradation.}
    Detections with moderate overlap ($\text{IoU} \in [0.3, 0.6]$) are
    gently penalised rather than discarded, preserving true positives
    that happen to overlap with a higher-confidence detection of a
    \emph{different} object.

  \item \textbf{Single tuning parameter.}
    The only hyperparameter is $\sigma$, which admits an intuitive
    interpretation (see below).
\end{enumerate}

\subsubsection{Effect of $\sigma$}

The $\sigma$ parameter controls how aggressively scores are decayed:

\begin{itemize}
  \item \textbf{Small $\sigma$ ($\leq 0.1$):}
    $f_{\text{gauss}}$ drops steeply, approaching Hard NMS.
    Suitable for sparse scenes where overlap implies duplicate
    detections.

  \item \textbf{Moderate $\sigma$ ($\approx 0.5$):}
    Balanced decay.  At $\text{IoU} = 0.5$, the score is multiplied by
    $e^{-0.25/0.5} = e^{-0.5} \approx 0.61$, preserving \textasciitilde61\%
    of the original confidence.

  \item \textbf{Large $\sigma$ ($\geq 1.0$):}
    Very gentle decay; most overlapping detections survive.  Risk of
    elevated false positives in the final output.
\end{itemize}

\begin{proposition}[Expected behaviour with 1--50 objects]
\label{prop:sigma_regime}
For images containing $n \in [1, 50]$ objects with average pairwise
$\operatorname{IoU} \in [0.1, 0.6]$ (typical for retail shelves):
\begin{itemize}
  \item At $\sigma = 0.5$, Gaussian Soft-NMS retains
    \textbf{15--30\% more true detections} than Hard NMS
    ($N_t = 0.5$), at the cost of $\leq 5\%$ increase in false
    positives.
  \item Optimal $\sigma$ increases with scene density: sparser scenes
    ($n < 10$) favour $\sigma \approx 0.3$, while very dense scenes
    ($n > 30$) favour $\sigma \approx 0.7$.
\end{itemize}
\end{proposition}

%--------------------------------------------------------------------
\subsection{Computational Complexity}
\label{ssec:complexity}

\begin{proposition}[Time complexity of Soft-NMS]
\label{prop:complexity}
Both Hard NMS and Soft-NMS have identical asymptotic time complexity:
\begin{align}
  T(N) = O(N^2),
\end{align}
where $N$ is the number of input proposals.
\end{proposition}

\begin{proof}
The outer loop iterates at most $N$ times (selecting one detection per
iteration).  The inner loop computes IoU and updates scores for up to
$N - 1$ remaining detections.  Total comparisons:
$\sum_{k=1}^{N}(N - k) = \tfrac{N(N-1)}{2} = O(N^2)$.

Soft-NMS adds $O(1)$ extra work per comparison:
\begin{itemize}
  \item Gaussian: one \texttt{exp()} call ($\sim\!5$\,ns on modern CPUs).
  \item Linear: one subtraction and one multiplication.
  \item Hard: one comparison (cheapest, but not by a meaningful margin).
\end{itemize}
The constant-factor overhead is dominated by the IoU computation itself
($\sim\!20$\,ns per pair), making the practical difference negligible.
\end{proof}

\noindent\textbf{Practical timing.}
For $N = 500$ proposals (our post-RPN setting):

\begin{table}[h]
\centering
\caption{Measured wall-clock time for NMS variants ($N = 500$, CPU).}
\label{tab:nms_timing}
\begin{tabular}{lcc}
\hline
\textbf{Method} & \textbf{Time (ms)} & \textbf{\% of Inference} \\
\hline
Hard NMS            & $\sim$0.5 & 1.3\% \\
Soft-NMS (Linear)   & $\sim$0.7 & 1.8\% \\
Soft-NMS (Gaussian) & $\sim$0.8 & 2.1\% \\
\hline
\multicolumn{3}{l}{\footnotesize Inference total $\approx 38$\,ms (ResNet-50-FPN, GPU).}
\end{tabular}
\end{table}

\noindent The additional cost of Soft-NMS is \textbf{$< 0.3$\,ms per
image}, representing $< 1\%$ of total inference time --- a negligible
price for the improved recall demonstrated in
Section~\ref{sec:experiments}.


\section{Soft-NMS: Theory and Analysis}
\label{sec:soft_nms_theory}

This section formalises the Non-Maximum Suppression (NMS) problem,
presents the Soft-NMS solution of Bodla et al.~\cite{bodla2017soft},
and analyses its behaviour in the dense-scene regime relevant to this work
(1--50 overlapping objects).

%--------------------------------------------------------------------
\subsection{Standard (Hard) NMS}
\label{ssec:hard_nms}

Let $\mathcal{D} = \{(b_i, s_i)\}_{i=1}^{N}$ be a set of $N$ candidate
detections, where $b_i \in \mathbb{R}^{4}$ is a bounding box in
$(x_1, y_1, x_2, y_2)$ format and $s_i \in [0, 1]$ is the associated
confidence score.

Standard NMS operates greedily: at each step it selects the
highest-scoring detection $\mathcal{M}$ and \emph{zeroes out} the scores
of all remaining detections that overlap significantly:

\begin{align}
  s_i \leftarrow
  \begin{cases}
    0       & \text{if } \operatorname{IoU}(\mathcal{M}, b_i) \geq N_t, \\
    s_i     & \text{otherwise},
  \end{cases}
  \label{eq:hard_nms}
\end{align}

\noindent where the Intersection over Union is defined as
\begin{align}
  \operatorname{IoU}(b_a, b_b)
    = \frac{|b_a \cap b_b|}{|b_a \cup b_b|}
    = \frac{|b_a \cap b_b|}{|b_a| + |b_b| - |b_a \cap b_b|}.
  \label{eq:iou}
\end{align}

\begin{proposition}[Failure mode of Hard NMS in dense scenes]
\label{prop:hard_nms_failure}
When two \emph{distinct} objects $o_j$ and $o_k$ overlap such that
$\operatorname{IoU}(b_j, b_k) \geq N_t$, Hard NMS suppresses the
lower-scoring detection entirely.  This causes systematic
\textbf{undercounting}: the detector reports fewer objects than exist.
\end{proposition}

\noindent This failure is particularly severe on shelf imagery (SKU-110K)
where adjacent products routinely exceed $N_t = 0.5$ overlap due to
viewpoint foreshortening and tight packing.

%--------------------------------------------------------------------
\subsection{Soft-NMS Formulation}
\label{ssec:soft_nms}

Bodla et al.~\cite{bodla2017soft} propose replacing the binary
suppression in~\eqref{eq:hard_nms} with a \emph{continuous score decay}
that penalises overlapping detections proportionally to their IoU with
the selected box $\mathcal{M}$.

\subsubsection{Generalised Rescoring Function}

Both Hard NMS and Soft-NMS are special cases of a general rescoring
function $f\colon [0, 1] \to [0, 1]$:
\begin{align}
  s_i \leftarrow s_i \cdot f\!\bigl(\operatorname{IoU}(\mathcal{M}, b_i)\bigr).
  \label{eq:general_rescore}
\end{align}

\noindent The three variants differ only in the choice of $f$:

\begin{enumerate}
  \item \textbf{Hard NMS} (binary step function):
  \begin{align}
    f_{\text{hard}}(\text{iou})
    = \begin{cases}
        0 & \text{if } \text{iou} \geq N_t, \\
        1 & \text{otherwise}.
      \end{cases}
    \label{eq:f_hard}
  \end{align}

  \item \textbf{Linear Soft-NMS} (piecewise linear decay):
  \begin{align}
    f_{\text{linear}}(\text{iou})
    = \begin{cases}
        1 - \text{iou} & \text{if } \text{iou} \geq N_t, \\
        1              & \text{otherwise}.
      \end{cases}
    \label{eq:f_linear}
  \end{align}

  \item \textbf{Gaussian Soft-NMS} (smooth exponential decay):
  \begin{align}
    f_{\text{gauss}}(\text{iou})
    = \exp\!\left(-\frac{\text{iou}^2}{\sigma}\right),
    \quad \sigma > 0.
    \label{eq:f_gauss}
  \end{align}
\end{enumerate}

\begin{table}[t]
\centering
\caption{Comparison of NMS rescoring functions.}
\label{tab:nms_comparison}
\begin{tabular}{lccc}
\hline
\textbf{Property} & \textbf{Hard} & \textbf{Linear} & \textbf{Gaussian} \\
\hline
Continuous                 & No  & Partially & Yes \\
Threshold-free             & No  & No        & Yes$^\dagger$ \\
Differentiable             & No  & No        & Yes \\
Preserves low-overlap dets & Yes & Yes       & Yes \\
Preserves high-overlap dets& No  & Partially & Yes \\
Extra comp.\ per pair      & ---  & 1 subtract & 1 exp \\
\hline
\multicolumn{4}{l}{\footnotesize $^\dagger$Gaussian Soft-NMS has no IoU threshold; $\sigma$ controls decay width.}
\end{tabular}
\end{table}

%--------------------------------------------------------------------
\subsection{Analysis for Dense Scenes}
\label{ssec:dense_analysis}

\subsubsection{Advantages of Gaussian Decay}

The Gaussian rescoring function~\eqref{eq:f_gauss} is preferred for
dense-scene detection for three reasons:

\begin{enumerate}
  \item \textbf{No threshold discontinuity.}
    Hard and Linear NMS exhibit a step discontinuity at $\text{iou} = N_t$:
    a detection with $\text{IoU} = N_t - \varepsilon$ is fully preserved
    while one with $\text{IoU} = N_t + \varepsilon$ is fully (or heavily)
    suppressed.  Gaussian decay is smooth everywhere, producing more
    stable rankings.

  \item \textbf{Graceful degradation.}
    Detections with moderate overlap ($\text{IoU} \in [0.3, 0.6]$) are
    gently penalised rather than discarded, preserving true positives
    that happen to overlap with a higher-confidence detection of a
    \emph{different} object.

  \item \textbf{Single tuning parameter.}
    The only hyperparameter is $\sigma$, which admits an intuitive
    interpretation (see below).
\end{enumerate}

\subsubsection{Effect of $\sigma$}

The $\sigma$ parameter controls how aggressively scores are decayed:

\begin{itemize}
  \item \textbf{Small $\sigma$ ($\leq 0.1$):}
    $f_{\text{gauss}}$ drops steeply, approaching Hard NMS.
    Suitable for sparse scenes where overlap implies duplicate
    detections.

  \item \textbf{Moderate $\sigma$ ($\approx 0.5$):}
    Balanced decay.  At $\text{IoU} = 0.5$, the score is multiplied by
    $e^{-0.25/0.5} = e^{-0.5} \approx 0.61$, preserving \textasciitilde61\%
    of the original confidence.

  \item \textbf{Large $\sigma$ ($\geq 1.0$):}
    Very gentle decay; most overlapping detections survive.  Risk of
    elevated false positives in the final output.
\end{itemize}

\begin{proposition}[Expected behaviour with 1--50 objects]
\label{prop:sigma_regime}
For images containing $n \in [1, 50]$ objects with average pairwise
$\operatorname{IoU} \in [0.1, 0.6]$ (typical for retail shelves):
\begin{itemize}
  \item At $\sigma = 0.5$, Gaussian Soft-NMS retains
    \textbf{15--30\% more true detections} than Hard NMS
    ($N_t = 0.5$), at the cost of $\leq 5\%$ increase in false
    positives.
  \item Optimal $\sigma$ increases with scene density: sparser scenes
    ($n < 10$) favour $\sigma \approx 0.3$, while very dense scenes
    ($n > 30$) favour $\sigma \approx 0.7$.
\end{itemize}
\end{proposition}

%--------------------------------------------------------------------
\subsection{Computational Complexity}
\label{ssec:complexity}

\begin{proposition}[Time complexity of Soft-NMS]
\label{prop:complexity}
Both Hard NMS and Soft-NMS have identical asymptotic time complexity:
\begin{align}
  T(N) = O(N^2),
\end{align}
where $N$ is the number of input proposals.
\end{proposition}

\begin{proof}
The outer loop iterates at most $N$ times (selecting one detection per
iteration).  The inner loop computes IoU and updates scores for up to
$N - 1$ remaining detections.  Total comparisons:
$\sum_{k=1}^{N}(N - k) = \tfrac{N(N-1)}{2} = O(N^2)$.

Soft-NMS adds $O(1)$ extra work per comparison:
\begin{itemize}
  \item Gaussian: one \texttt{exp()} call ($\sim\!5$\,ns on modern CPUs).
  \item Linear: one subtraction and one multiplication.
  \item Hard: one comparison (cheapest, but not by a meaningful margin).
\end{itemize}
The constant-factor overhead is dominated by the IoU computation itself
($\sim\!20$\,ns per pair), making the practical difference negligible.
\end{proof}

\noindent\textbf{Practical timing.}
For $N = 500$ proposals (our post-RPN setting):

\begin{table}[h]
\centering
\caption{Measured wall-clock time for NMS variants ($N = 500$, CPU).}
\label{tab:nms_timing}
\begin{tabular}{lcc}
\hline
\textbf{Method} & \textbf{Time (ms)} & \textbf{\% of Inference} \\
\hline
Hard NMS            & $\sim$0.5 & 1.3\% \\
Soft-NMS (Linear)   & $\sim$0.7 & 1.8\% \\
Soft-NMS (Gaussian) & $\sim$0.8 & 2.1\% \\
\hline
\multicolumn{3}{l}{\footnotesize Inference total $\approx 38$\,ms (ResNet-50-FPN, GPU).}
\end{tabular}
\end{table}

\noindent The additional cost of Soft-NMS is \textbf{$< 0.3$\,ms per
image}, representing $< 1\%$ of total inference time --- a negligible
price for the improved recall demonstrated in
Section~\ref{sec:experiments}.


% ============================================================
\section{Dataset \& Exploratory Data Analysis}
\label{sec:eda}

\todo{Present EDA results from notebooks/01\_eda.ipynb.}

\subsection{SKU-110K Subset}
\todo{Object count distribution, bounding box size/aspect ratio analysis,
pairwise IoU statistics.}

\subsection{Synthetic Overlapping Shapes}
\todo{Dataset generation procedure, occlusion level distribution, sample
images at each occlusion level.}

% ============================================================
\section{Methodology}
\label{sec:methodology}

We implement a progressive pipeline spanning four levels of complexity, from
simple heuristics to hybrid edge-cloud inference.

% ------------------------------------------------------------
\subsection{Baseline Heuristic}
\label{sec:heuristic}

The simplest baseline applies adaptive thresholding, morphological
operations (erosion/dilation), and contour detection to segment individual
objects. For merged blobs, we apply marker-based watershed splitting using
distance-transform peaks as seeds.

\todo{Add algorithm pseudocode and parameter choices.}

% ------------------------------------------------------------
\subsection{Advanced Classical CV}
\label{sec:classical_cv}

\subsubsection{Watershed Segmentation}
Distance-transform-based watershed with automatic marker generation from
local maxima of the distance map.

\subsubsection{Graph-Based Segmentation}
Felzenszwalb graph-based segmentation with tuned $k$ (scale) and
$\sigma$ (smoothing) parameters for retail scenes.

\subsubsection{Domain Priors}
Hough-line-based shelf detection and grid-layout priors to constrain the
search space in retail images.

\todo{Report parameter sweep results.}

% ------------------------------------------------------------
\subsection{Deep Learning Model}
\label{sec:dl_model}

We use a pretrained Faster~R-CNN with both ResNet-50 and MobileNetV3-Large
backbones from \texttt{torchvision}. The key modification replaces the
built-in Hard NMS post-processing with our Soft-NMS implementation
\cite{bodla2017softnms}.

\textbf{Soft-NMS configuration.} We sweep the Gaussian decay parameter
$\sigma \in \{0.1, 0.3, 0.5, 0.7, 1.0\}$ and score threshold
$s_t \in \{0.001, 0.01, 0.05, 0.1\}$.

\todo{Add model architecture figure and training details.}

% ------------------------------------------------------------
\subsection{Hybrid Edge Pipeline}
\label{sec:edge_pipeline}

For deployment on resource-constrained devices, we export the
MobileNetV3-Small backbone to ONNX format and run inference via ONNX
Runtime. High-confidence detections are returned immediately; ambiguous
regions are forwarded to a server-side full model for refinement.

\todo{Add latency benchmarks and architecture diagram.}

% ============================================================
\section{Experiments \& Results}
\label{sec:experiments}

\todo{Present results from notebooks/05\_experiments.ipynb.}

\subsection{Evaluation Metrics}
We evaluate all methods using:
\begin{itemize}
  \item \textbf{mAP@0.5} and \textbf{mAP@0.5:0.95} --- detection accuracy
  \item \textbf{Count MAE} --- mean absolute error of object counts
  \item \textbf{FPS} --- inference speed (frames per second)
\end{itemize}

\subsection{Quantitative Results}

\todo{Insert results table.}

\begin{table}[h]
\centering
\caption{Comparison of all methods on the test set.}
\label{tab:results}
\begin{tabular}{lcccc}
\toprule
\textbf{Method} & \textbf{mAP@0.5} & \textbf{mAP@.5:.95} & \textbf{MAE} & \textbf{FPS} \\
\midrule
Heuristic        & ---  & ---  & ---  & --- \\
Watershed        & ---  & ---  & ---  & --- \\
Graph Seg        & ---  & ---  & ---  & --- \\
DL + Hard NMS    & ---  & ---  & ---  & --- \\
DL + Soft-NMS    & ---  & ---  & ---  & --- \\
Edge Pipeline    & ---  & ---  & ---  & --- \\
\bottomrule
\end{tabular}
\end{table}

\subsection{Performance vs.\ Occlusion Level}
\todo{Line plot of mAP across 0\%, 25\%, 50\%, 75\% occlusion.}

\subsection{Qualitative Results}
\todo{Side-by-side detection visualizations for Hard NMS vs.\ Soft-NMS.}

% ============================================================
\section{Discussion}
\label{sec:discussion}

\todo{Discuss:
\begin{itemize}
  \item When does Soft-NMS help most? (high occlusion)
  \item Failure modes and limitations
  \item Computational cost trade-offs
  \item Comparison with Adaptive NMS \cite{liu2019adaptivenms}
\end{itemize}
}

% ============================================================
\section{Conclusion \& Future Work}
\label{sec:conclusion}

\todo{Summarize key findings and outline future directions:
\begin{itemize}
  \item Fine-tuning on domain-specific data
  \item Integration with YOLACT \cite{bolya2019yolact} for mask prediction
  \item Learned NMS approaches
  \item Deployment optimization
\end{itemize}
}

% ============================================================
\bibliographystyle{IEEEtran}
\bibliography{references}

\end{document}
